\documentclass[12pt,a4paper]{article}
\usepackage{amsmath,amssymb}
\usepackage[UTF8]{ctex}
\usepackage[margin=2.5cm]{geometry}
\usepackage{booktabs}
\usepackage{tabularx}

\title{问题1 公式与推导过程}
\date{}

\begin{document}
\maketitle

%=============================================================================
% 问题1.1
%=============================================================================
\section{问题1.1：未来五年老人数量预测}

\subsection{问题重述}

已知各小区当前60岁以上老人的自理、半失能、失能三类人数，以及年死亡率、年新增率、状态转移概率等参数，要求预测未来五年各小区三类老人的人数变化，并计算第1年总人口增长率和5年复合年均增长率。

\subsection{问题分析}

本问的核心是建立一个人口动态演化模型。老人总人口受两个因素影响：新增人口（按年新增率）和死亡人口（按年死亡率）。同时，三类老人之间存在单向状态转移：自理$\to$半失能$\to$失能。为提高预测精度，采用逐日递推的方式，将年度参数转化为日度参数后进行迭代计算。

\subsection{模型假设}

\begin{enumerate}
    \item 将老人群体视为封闭系统，人口变化仅由新增和死亡决定，不考虑迁入迁出。
    \item 每天新增的老人全部归入自理类，新增时不存在半失能或失能状态。
    \item 年死亡率 $M$ 对所有老人和所有小区均相同，不随年龄、健康状态变化。
    \item 状态转移是不可逆的单向过程，即自理$\to$半失能$\to$失能，不存在反向恢复。
    \item 每天的操作顺序为：新增$\to$死亡$\to$状态转移，新增人口与存量老人并行面对当日死亡率。
    \item 每天的存活事件、转移事件相互独立，可使用乘法模型计算累积效果。
    \item 因日转移率 $t_1, t_2 \ll 1$，单日转移人数近似为当前人数乘以日转移率（线性近似）。
    \item 年死亡率 $M$、年新增率 $G$、年转移概率 $p_1, p_2$ 在5年预测期内保持不变。
\end{enumerate}

\subsection{符号说明}

本文所用符号及其含义如表~\ref{tab:symbols_1.1}~所示。

\begin{table}[htbp]
\centering
\caption{问题1.1 符号说明}
\label{tab:symbols_1.1}
\begin{tabular}{lll}
\toprule
符号 & 含义 & 单位/取值范围 \\
\midrule
$C_t$ & 第 $t$ 天自理老人人数 & 人，$C_t \geq 0$ \\
$S_t$ & 第 $t$ 天半失能老人人数 & 人，$S_t \geq 0$ \\
$D_t$ & 第 $t$ 天失能老人人数 & 人，$D_t \geq 0$ \\
$T_t$ & 第 $t$ 天老人总人数，$T_t = C_t + S_t + D_t$ & 人 \\
$C_0, S_0, D_0$ & 初始时刻（第0天）各小区的三类老人人数 & 人 \\
$M$ & 年死亡率，题目给定 $M=0.05$ & 无量纲，$0 < M < 1$ \\
$G$ & 年新增率，题目给定 $G=0.07$ & 无量纲，$G > 0$ \\
$p_1$ & 自理$\to$半失能的年转移概率 & 无量纲，$0 < p_1 < 1$ \\
$p_2$ & 半失能$\to$失能的年转移概率 & 无量纲，$0 < p_2 < 1$ \\
$\mu$ & 日死亡率（待推导） & 无量纲 \\
$t_1$ & 自理$\to$半失能的日转移率（待推导） & 无量纲 \\
$t_2$ & 半失能$\to$失能的日转移率（待推导） & 无量纲 \\
\bottomrule
\end{tabular}
\end{table}

\subsection{模型建立}

\subsubsection{日度参数转换}

模型采用逐日递推，需将年度参数转化为日度参数。核心思想：若一年365天每天施加相同的日度率，则累积效果应等于年度率。

\noindent\textbf{日死亡率 $\mu$ 的推导}

已知年死亡率 $M=0.05$，即一个老人存活一年后死亡的概率为 $M$，存活概率为 $1-M$。

设日死亡率为 $\mu$，则任意一天存活的概率为 $1-\mu$。由于各天存活事件相互独立，连续存活365天的概率为 $(1-\mu)^{365}$。根据年存活概率的定义，可得：

\begin{equation}
(1-\mu)^{365} = 1 - M
\tag{1.1}
\end{equation}

对式~(1.1)两边取365次方根，解得日死亡率：

\begin{equation}
\mu = 1 - (1 - M)^{1/365}
\tag{1.2}
\end{equation}

其中：$\mu$ 为日死亡率（无量纲），$M$ 为年死亡率，题目给定 $M=0.05$。

代入 $M=0.05$，得 $\mu \approx 0.0001405$。

\noindent\textbf{日转移率 $t_1$、$t_2$ 的推导}

年转移概率 $p_1$ 表示一个自理老人在一年内转为半失能的概率。采用与死亡率相同的衰减逻辑：每天有一定的转移率 $t_1$，则一年内未转移的概率为 $(1-t_1)^{365}$，应等于 $1-p_1$：

\begin{equation}
(1-t_1)^{365} = 1 - p_1
\tag{1.3}
\end{equation}

对式~(1.3)两边取365次方根，解得自理$\to$半失能的日转移率：

\begin{equation}
t_1 = 1 - (1 - p_1)^{1/365}
\tag{1.4}
\end{equation}

其中：$t_1$ 为自理$\to$半失能的日转移率（无量纲），$p_1$ 为对应的年转移概率。

同理，对半失能$\to$失能的转移，可得：

\begin{equation}
(1-t_2)^{365} = 1 - p_2
\tag{1.5}
\end{equation}

\begin{equation}
t_2 = 1 - (1 - p_2)^{1/365}
\tag{1.6}
\end{equation}

其中：$t_2$ 为半失能$\to$失能的日转移率（无量纲），$p_2$ 为对应的年转移概率。

\noindent\textbf{年新增方式}

年新增率 $G=0.07$ 表示每年初在上年末总人口的基础上增加 $G$ 倍，新增人口全部归入自理类。即第 $k$ 年初的实际执行方式为：

\begin{equation}
C_{k}^{(0)} = C_{k-1} + T_{k-1} \cdot G
\tag{1.7}
\end{equation}

其中：$C_{k}^{(0)}$ 为第 $k$ 年初新增后的自理人数（单位：人），$C_{k-1}$ 为第 $k-1$ 年末自理人数（单位：人），$T_{k-1}=C_{k-1}+S_{k-1}+D_{k-1}$ 为第 $k-1$ 年末总人数（单位：人），$G$ 为年新增率，$G=0.07$。

\subsubsection{年-日双层递推模型}

模型采用年-日双层结构：每年初新增一次，之后以日为步长循环365次，依次执行死亡和状态转移。这一时序安排的理由是：新增人口在年初一次性加入后，与存量老人同样在全年各日面对死亡率和状态退化风险。

\noindent\textbf{外层循环——年初新增}

第 $k$ 年初，在上年末总人口 $T_{k-1}$ 的基础上按年新增率 $G$ 增加人口，全部归入自理类：

\begin{equation}
\left\{
\begin{aligned}
C_k^{(0)} &= C_{k-1} + T_{k-1} \cdot G \\
S_k^{(0)} &= S_{k-1} \\
D_k^{(0)} &= D_{k-1}
\end{aligned}
\right.
\tag{1.8}
\end{equation}

其中：$C_k^{(0)}, S_k^{(0)}, D_k^{(0)}$ 为第 $k$ 年逐日递推的初始值（单位：人），$C_{k-1}, S_{k-1}, D_{k-1}$ 为第 $k-1$ 年末的三类人数（单位：人），$T_{k-1} = C_{k-1} + S_{k-1} + D_{k-1}$ 为第 $k-1$ 年末总人数（单位：人），$G$ 为年新增率（无量纲），$G=0.07$。

\noindent\textbf{内层循环——逐日死亡与转移}

以第 $k$ 年初始值 $C_k^{(0)}, S_k^{(0)}, D_k^{(0)}$ 为起点，按天迭代 $365$ 次，每天依次执行死亡和状态转移两个阶段。

\noindent\textbf{Step 1：死亡（三类同步衰减）}

每天所有老人均面对日死亡率 $\mu$，三类人数同时乘以存活率 $(1-\mu)$：

\begin{equation}
\left\{
\begin{aligned}
C' &= C_{t-1} \cdot (1 - \mu) \\
S' &= S_{t-1} \cdot (1 - \mu) \\
D' &= D_{t-1} \cdot (1 - \mu)
\end{aligned}
\right.
\tag{1.9}
\end{equation}

其中：$C', S', D'$ 为死亡筛选后的三类人数（中间变量，单位：人），$C_{t-1}, S_{t-1}, D_{t-1}$ 为前一天的三类人数（单位：人），$\mu$ 为日死亡率（无量纲）。

\noindent\textbf{Step 2：状态转移}

死亡筛选后，自理老人以日转移率 $t_1$ 转为半失能，半失能老人以日转移率 $t_2$ 转为失能。由于 $t_1, t_2 \ll 1$，单日转移人数可近似为当前人数乘以日转移率（线性近似）：

\begin{equation}
\left\{
\begin{aligned}
\Delta_C &= C' \cdot t_1 \\
\Delta_S &= S' \cdot t_2
\end{aligned}
\right.
\tag{1.10}
\end{equation}

其中：$\Delta_C$ 为当日自理转半失能的人数（单位：人），$\Delta_S$ 为当日半失能转失能的人数（单位：人），$t_1, t_2$ 为对应的日转移率（无量纲）。

更新后的三类人数为：

\begin{equation}
\left\{
\begin{aligned}
C_t &= C' - \Delta_C = C'(1 - t_1) \\
S_t &= S' + \Delta_C - \Delta_S = S' + C' \cdot t_1 - S' \cdot t_2 \\
D_t &= D' + \Delta_S = D' + S' \cdot t_2
\end{aligned}
\right.
\tag{1.11}
\end{equation}

其中：$C_t, S_t, D_t$ 为第 $t$ 天结束时的人数（单位：人）。

\noindent\textbf{逐日递推的完整形式}

将式~(1.9)、(1.10)、(1.11)合并，消去中间变量，得到从第 $t-1$ 天到第 $t$ 天的完整递推公式：

\begin{equation}
\boxed{
\left\{
\begin{aligned}
C_t &= C_{t-1} (1 - \mu) (1 - t_1) \\
S_t &= S_{t-1} (1 - \mu) (1 - t_2) + C_{t-1} (1 - \mu) \cdot t_1 \\
D_t &= D_{t-1} (1 - \mu) + S_{t-1} (1 - \mu) \cdot t_2
\end{aligned}
\right.
}
\tag{1.12}
\end{equation}

其中：$C_{t-1}, S_{t-1}, D_{t-1}$ 为前一天的三类人数（单位：人），$\mu$ 为日死亡率（无量纲），$t_1, t_2$ 为日转移率（无量纲）。

将式~(1.12)逐日迭代 $365$ 次，得到第 $k$ 年的年末人数：

\begin{equation}
(C_k, S_k, D_k) = \text{迭代365次后的} \; (C_t, S_t, D_t)
\tag{1.13}
\end{equation}

\noindent\textbf{完整算法流程}

\begin{enumerate}
    \item 令 $(C_0, S_0, D_0)$ 为第0年末（即第1年初）各小区实际人数。
    \item 对每个年份 $k = 1, 2, \dots, 5$：
    \begin{enumerate}
        \item 按式~(1.8)执行年初新增，得到 $C_k^{(0)}, S_k^{(0)}, D_k^{(0)}$。
        \item 以 $C_k^{(0)}, S_k^{(0)}, D_k^{(0)}$ 为初始值，将式~(1.12)迭代 $365$ 次，得到年末人数 $C_k, S_k, D_k$。
    \end{enumerate}
    \item 输出各年末的三类人口数据。
\end{enumerate}

\subsection{增长率指标}

\subsubsection{第1年增长率}

第1年结束后总人口为 $T_{365}$，初始为 $T_0$，则第1年总人口增长率为：

\begin{equation}
g_1 = \frac{T_{365} - T_0}{T_0}
\tag{1.14}
\end{equation}

其中：$g_1$ 为第1年的总人口增长率（无量纲），$T_{365}$ 为第1年末总人口（单位：人），$T_0$ 为初始总人口（单位：人）。

\subsubsection{5年复合年均增长率}

设第5年末总人口为 $T_{1825}$（即递推1825天后的总人口）。复合年均增长率（CAGR）定义为：假设每年以恒定速率增长，5年后从 $T_0$ 增长到 $T_{1825}$，即满足 $T_0(1+r)^5 = T_{1825}$，解得：

\begin{equation}
r = \left(\frac{T_{1825}}{T_0}\right)^{1/5} - 1
\tag{1.15}
\end{equation}

其中：$r$ 为5年复合年均增长率（无量纲），$T_{1825}$ 为第5年末总人口（单位：人），$T_0$ 为初始总人口（单位：人）。


%=============================================================================
% 问题1.2
%=============================================================================
\section{问题1.2：第5年末各小区理论月需求计算}

\subsection{问题重述}

已知各服务项目（助餐、日间照料、上门护理、康复理疗、助浴、紧急救助）的月均需求次数，以及问题1.1预测得到的第5年末各小区三类老人人数，要求计算第5年末各小区对各服务项目的理论月需求次数。

\subsection{问题分析}

本问的需求计算基于一个简单的线性模型：需求次数与人数成正比。对于每类老人，其对某项服务的月需求次数等于人数乘以人均需求率。由于需求次数必须为整数，计算结果需进行四舍五入取整。最终将三类老人的需求汇总，得到各小区对各服务项目的总需求。

\subsection{模型假设}

\begin{enumerate}
    \item 各服务的月需求次数率（人均需求率）在5年内保持不变。
    \item 需求次数与人数成正比，即需求率不随人数规模变化。
    \item 不同老人之间的需求相互独立，可直接累加。
    \item 需求次数必须为整数，采用四舍五入取整。
\end{enumerate}

\subsection{符号说明}

本文所用符号及其含义如表~\ref{tab:symbols_1.2}~所示。

\begin{table}[htbp]
\centering
\caption{问题1.2 符号说明}
\label{tab:symbols_1.2}
\begin{tabular}{@{} l p{4.5cm} p{3.5cm} @{}}
\toprule
符号 & 含义 & 单位/取值范围 \\
\midrule
$i$ & 小区编号 & $i = A, B, C, \ldots, J$ \\
$n_C^{(i)}, n_S^{(i)}, n_D^{(i)}$ & 小区 $i$ 第5年末三类老人人数（即 $n_C^{(i)} = C_{1825}^{(i)}$ 等） & 人，$n \geq 0$ \\
$svc$ & 服务项目 & 助餐、日间照料、上门护理、康复理疗、助浴、紧急救助 \\
$r_{svc}^C, r_{svc}^S, r_{svc}^D$ & 三类老人每人每月对服务 $svc$ 的需求次数 & 次/人/月 \\
$\tilde{q}_{i,svc}^C$ & 小区 $i$ 自理老人对服务 $svc$ 的理论月需求次数（取整前） & 次 \\
$q_{i,svc}^C$ & 小区 $i$ 自理老人对服务 $svc$ 的理论月需求次数（取整后） & 次，整数 \\
$Q_{i,svc}$ & 小区 $i$ 对服务 $svc$ 的三类老人理论月需求总和 & 次，整数 \\
\bottomrule
\end{tabular}
\end{table}

\subsection{模型建立}

\subsubsection{单类老人需求计算}

对于小区 $i$ 的某类老人（以自理为例），第5年末人数为 $n_C^{(i)}$，该类老人每人每月对服务 $svc$ 的需求次数为 $r_{svc}^C$。则小区 $i$ 自理老人对服务 $svc$ 的理论月需求次数（取整前）为人数乘以人均需求率：

\begin{equation}
\tilde{q}_{i,svc}^C = n_C^{(i)} \cdot r_{svc}^C
\tag{2.1}
\end{equation}

其中：$\tilde{q}_{i,svc}^C$ 为取整前的理论需求次数（单位：次），$n_C^{(i)}$ 为小区 $i$ 第5年末自理老人人数（单位：人），$r_{svc}^C$ 为自理老人对服务 $svc$ 的人均月需求次数（单位：次/人/月）。

由于需求次数必须为整数，对式~(2.1)取四舍五入：

\begin{equation}
q_{i,svc}^C = \text{round}\bigl(n_C^{(i)} \cdot r_{svc}^C\bigr)
\tag{2.2}
\end{equation}

其中：$\text{round}(\cdot)$ 表示四舍五入取整函数，$q_{i,svc}^C$ 为取整后的实际需求次数（单位：次，整数）。

同理，半失能和失能老人的需求计算公式为：

\begin{equation}
\left\{
\begin{aligned}
q_{i,svc}^S &= \text{round}\bigl(n_S^{(i)} \cdot r_{svc}^S\bigr) \\
q_{i,svc}^D &= \text{round}\bigl(n_D^{(i)} \cdot r_{svc}^D\bigr)
\end{aligned}
\right.
\tag{2.3}
\end{equation}

\subsubsection{总需求计算}

小区 $i$ 对服务 $svc$ 的总需求为三类老人需求之和：

\begin{equation}
Q_{i,svc} = q_{i,svc}^C + q_{i,svc}^S + q_{i,svc}^D
\tag{2.4}
\end{equation}

其中：$Q_{i,svc}$ 为小区 $i$ 对服务 $svc$ 的理论月需求总和（单位：次，整数），$q_{i,svc}^C, q_{i,svc}^S, q_{i,svc}^D$ 分别为自理、半失能、失能老人的需求次数（单位：次）。


%=============================================================================
% 问题1.3
%=============================================================================
\section{问题1.3：消费约束下实际月需求计算}

\subsection{问题重述}

已知各小区人均月收入、各服务项目营收单价、三类老人的消费比例上限，以及问题1.2计算得到的理论月需求次数，要求在消费约束条件下计算各小区各服务项目的实际月需求次数。当理论消费超出上限时，采用等比例缩减策略。

\subsection{问题分析}

本问引入消费约束，即每类老人每月用于养老服务的消费不超过月收入的一定比例。当理论消费超出上限时，需要对需求进行缩减。缩减策略采用等比例方式，即所有服务项目的需求次数按相同比例缩小，以保证缩减的公平性。

\subsection{模型假设}

\begin{enumerate}
    \item 消费上限是硬约束，超出上限时必须缩减需求。
    \item 当需要缩减时，所有服务项目的需求次数按相同比例缩小，不偏袒任何单项服务。
    \item 同一小区同一类老人的缩减系数对所有服务项目相同。
    \item 各小区人均月收入 $I_i$ 在5年内保持不变。
    \item 各服务营收单价 $p_{svc}$ 在5年内保持不变。
    \item 每类老人的消费决策相互独立，不存在跨类别的消费转移。
\end{enumerate}

\subsection{符号说明}

本文所用符号及其含义如表~\ref{tab:symbols_1.3}~所示。

\begin{table}[htbp]
\centering
\caption{问题1.3 符号说明}
\label{tab:symbols_1.3}
\begin{tabular}{@{} l p{5.5cm} l @{}}
\toprule
符号 & 含义 & 单位/取值范围 \\
\midrule
$I_i$ & 小区 $i$ 人均月收入 & 元，$I_i > 0$ \\
$k_C, k_S, k_D$ & 自理、半失能、失能老人的消费比例上限 & 无量纲，$0 < k < 1$ \\
$p_{svc}$ & 服务 $svc$ 的营收单价 & 元/次，$p_{svc} > 0$ \\
$L_{i,C}, L_{i,S}, L_{i,D}$ & 小区 $i$ 自理、半失能、失能老人每月消费上限 & 元 \\
$w_{i,C}, w_{i,S}, w_{i,D}$ & 小区 $i$ 自理、半失能、失能老人人均理论月消费 & 元 \\
$f_{i,C}, f_{i,S}, f_{i,D}$ & 小区 $i$ 自理、半失能、失能老人的需求缩减系数 & 无量纲，$0 < f \leq 1$ \\
$a_{i,svc}^C, a_{i,svc}^S, a_{i,svc}^D$ & 小区 $i$ 自理、半失能、失能老人对服务 $svc$ 的实际月需求次数 & 次，整数 \\
\bottomrule
\end{tabular}
\end{table}

\subsection{模型建立}

\subsubsection{消费上限计算}

每类老人每月用于养老服务的消费不超过月收入的一定比例。该比例由题目给定，按自理、半失能、失能分别设定上限系数。小区 $i$ 各类老人的消费上限为：

\begin{equation}
\left\{
\begin{aligned}
L_{i,C} &= I_i \cdot k_C \\
L_{i,S} &= I_i \cdot k_S \\
L_{i,D} &= I_i \cdot k_D
\end{aligned}
\right.
\tag{3.1}
\end{equation}

其中：$L_{i,C}, L_{i,S}, L_{i,D}$ 为小区 $i$ 自理、半失能、失能老人每月消费上限（单位：元），$I_i$ 为小区 $i$ 人均月收入（单位：元），$k_C, k_S, k_D$ 为对应的消费比例上限（无量纲）。

\subsubsection{人均理论月消费计算}

一个类型为 $C$（自理）的老人，每月对各服务的需求次数为 $r_{svc}^C$，每项服务单价为 $p_{svc}$。若不考虑消费约束，其月消费为所有服务的需求次数乘以对应单价之和：

\begin{equation}
w_{i,C} = \sum_{svc} r_{svc}^C \cdot p_{svc}
\tag{3.2}
\end{equation}

其中：$w_{i,C}$ 为自理老人人均理论月消费（单位：元），$r_{svc}^C$ 为自理老人对服务 $svc$ 的人均月需求次数（单位：次/人/月），$p_{svc}$ 为服务 $svc$ 的营收单价（单位：元/次）。

$w_{i,S}$、$w_{i,D}$ 的计算同理。

\subsubsection{缩减系数的推导}

判断是否需要缩减：比较人均理论月消费 $w_{i,C}$ 与消费上限 $L_{i,C}$。

若 $w_{i,C} \leq L_{i,C}$，理论消费在上限范围内，不需要缩减，系数 $f_{i,C}=1$。

若 $w_{i,C} > L_{i,C}$，理论消费超出上限。采用等比例缩减策略：所有服务项目的需求次数按相同比例缩小，使得缩减后的人均消费恰好等于上限。

以自理老人为例，设缩减系数为 $f$（$0 < f < 1$），缩减后每项服务的需求次数变为 $r_{svc}^C \cdot f$，则缩减后的人均消费为：

\begin{equation}
\sum_{svc} (r_{svc}^C \cdot f) \cdot p_{svc}
= f \cdot \sum_{svc} r_{svc}^C \cdot p_{svc}
= f \cdot w_{i,C}
\tag{3.3}
\end{equation}

令式~(3.3)等于消费上限 $L_{i,C}$，即 $f \cdot w_{i,C} = L_{i,C}$，解得缩减系数：

\begin{equation}
f = \frac{L_{i,C}}{w_{i,C}}
\tag{3.4}
\end{equation}

合并两种情况（不需缩减和需要缩减），用 $\min$ 函数统一表达：

\begin{equation}
f_{i,C} = \min\left(1,\; \frac{L_{i,C}}{w_{i,C}}\right)
\tag{3.5}
\end{equation}

其中：$f_{i,C}$ 为小区 $i$ 自理老人的需求缩减系数（无量纲，$0 < f_{i,C} \leq 1$），$L_{i,C}$ 为消费上限（单位：元），$w_{i,C}$ 为人均理论月消费（单位：元）。$f_{i,S}$、$f_{i,D}$ 的计算同理。

\subsubsection{实际需求计算}

将式~(2.1)的理论需求乘以缩减系数~(3.5)后再取整，得到实际需求次数：

\begin{equation}
a_{i,svc}^C = \text{round}\bigl(n_C^{(i)} \cdot r_{svc}^C \cdot f_{i,C}\bigr)
\tag{3.6}
\end{equation}

其中：$a_{i,svc}^C$ 为小区 $i$ 自理老人对服务 $svc$ 的实际月需求次数（单位：次，整数），$n_C^{(i)}$ 为第5年末自理人数（单位：人），$r_{svc}^C$ 为自理老人对服务 $svc$ 的人均月需求次数（单位：次/人/月），$f_{i,C}$ 为缩减系数（无量纲）。

$a_{i,svc}^S$、$a_{i,svc}^D$ 的计算同理。

\subsubsection{等比例缩减的性质}

由于缩减系数 $f_{i,C}$ 对同一小区同一类老人的所有服务项目相同，因此缩减后各服务之间的需求比例保持不变：

\begin{equation}
\frac{a_{i,svc_1}^C}{a_{i,svc_2}^C}
\approx \frac{r_{svc_1}^C}{r_{svc_2}^C}
\tag{3.7}
\end{equation}

这保证了缩减的公平性——不偏袒任何单项服务。

\end{document}
